% Evans second edition, section 2.6.
% English source: printed p. 86; PDF page 103 (1-based).
% Chinese section starts checked from chapter opening, PDF 31; mean-value pages checked at PDF 36-37 and 53-55.
% Status: editorial study draft with reconstructed proofs and mathematical annotations.
\section{References}
\label{sec:2-6}

The primary source is Lawrence C. Evans, \emph{Partial Differential
Equations}, second edition, Chapter 2, printed pp.~17--87
(local PDF pp.~34--104). The exposition here is reconstructed for study;
source theorem numbering is retained only as an explicit source identifier.

Evans's \S2.6 attributes the heat-ball theorem and the associated
strong maximum principle to N.~Watson, \emph{Proceedings of the London
Mathematical Society} 26 (1973), pp.~385--417.
The proof above follows the source's scaling/logarithm method and adds
the full mass computation, tip estimates, scale-defect identity, and
moment interpretation. The latter supplements are editorial deductions;
no separate Watson-paper comparison is claimed.
The English source omits the normalization calculation; the Chinese
translator's footnote on printed p.~37 suggests polar coordinates and
a Gamma substitution. The present proof spells out those constants.

The source points to Gilbarg--Trudinger, Chapters 2--4, for elliptic
theory; John and Friedman for the heat equation; Mikhailov for
regularity; Payne for backward uniqueness; and Folland and Strauss
for wave representation formulas. These are source reading directions,
not claims that those books were independently consulted for this draft.

The Chinese chapter opening (PDF p.~31) provides section starts at
printed pp.~13, 14, 30, 43, 56, and 60 for \S2.1--\S2.6.
Its mean-value and heat-ball pages were visually checked for terminology
and formulas as recorded in the chapter log. This is not a claim of
a line-by-line comparison of the entire Chinese chapter.

\paragraph{Scope and review.}
Elementary calculus, measure-theoretic limit theorems, the Gaussian
integral and Gamma integral, and the explicitly invoked complex-analytic
parameter-integral/Cauchy facts are prerequisites.
General-domain existence of harmonic boundary correctors is left to
later Dirichlet theory. All four model equations, their representation
methods, and the 24 source exercises have study explanations here;
mathematical self-checks and numerical moment checks are not independent
mathematical review.
